\section{Lean Development and Certification Scope}\label{app-sec:lean-formalization}
The repository contains a substantial Lean development of finite-dimensional inequalities, quotient geometry, scalar convergence recurrences, and conditional convergence interfaces. Organizing these reusable results independently of the article is useful: KL inequalities and recurrence lemmas, for example, have applications beyond transport. The entry point is \texttt{FlowSinkhorn.KLProjection}; thematic modules group analysis, bounds, geometry, applications, and certification infrastructure. \texttt{StatementMap.lean} provides stable aliases to historical paper-facing endpoints rather than new proofs.

\paragraph{What a successful build certifies.}
Lean checks a theorem under its declared assumptions. A theorem that assumes a bias estimate, a mass bound, an objective-gap inequality, or an already constructed certificate does not prove those ingredients. An alias to such a theorem, a generated comparison interface, and an absence of \texttt{sorry} do not establish that the article's original variational assumptions imply its conclusion. In particular, compiling the current project is not an end-to-end certificate for every statement of this revised article. The revised bounded-mass lemma, initialization, normalized graph constants, quotient pairing distinction, and graph scalar-order argument require statement-level synchronization and, in some cases, additional formal proofs.

\paragraph{Current status and reproduction.}
The revision audit is recorded in \texttt{lean/PAPER\_COVERAGE.md}. Historical statement aliases are retained to avoid breaking downstream imports, but their numerical labels must not be read as the revised paper's numbering. Running \texttt{lake build} in \texttt{lean/} checks the available Lean sources against the pinned toolchain and dependencies. It does not check this PDF, scientific relevance of assumptions, or external numerical experiments. The CI build has the same scope. The project is therefore best described as a partially formalized companion library with conditional interfaces; completion requires independent comparison of each Lean statement's hypotheses and conclusion with its mathematical counterpart, followed by formalization of every missing bridge. No percentage or blanket ``fully formalized'' claim is made here.
